%% notes-tex/032-amino-acid-module/sections/4-concordance.tex
\section{The second amino-acid screen\secstatus{todo}}
\label{sec:concordance}

The second amino-acid screen does not corroborate the first
(Table~\ref{tab:amino-acid-concordance}; \texttt{amino\_acid\_dataset\_concordance.py}).
Cooper 2010 profiled the auxotrophic MATa collection after overnight growth in synthetic
complete medium by capillary electrophoresis, 17 peaks each a linear ratio to its plate
mean, four of them composites and lysine split in two; M\"ulleder 2016 measured 19 amino
acids in mM on the prototrophic collection on minimal medium. Over the 4,051 strains both
screens hold, the matched pairs correlate at $-0.006$ (alanine) to $+0.059$ (glycine,
$p$ $2\times10^{-4}$), with glutamate at 0.043 and aspartate at 0.031 the only others
off zero; the top 5 percent of strains overlap at chance on every pair (the largest
enrichment 2.2, lysine peak b, 22 of 202 against 10 expected); and across all 17 by 19
cells the median $|\rho|$ is 0.018 and the maximum 0.127, on an unmatched pair, so the
matched cells do not stand out from the background. The two retest references bound
this: M\"ulleder's 194 strains measured twice give a mean Spearman of 0.20 (Pearson on
$\log_{10}$ 0.43, the 027 read), and Cooper's 47 duplicated genes, paired from the
loader's ledger against the served rows, give 0.06 (log Pearson 0.10). The Cooper pairs
are the screen's own within-study replicate agreement: the 47 genes the collection
carries as two deletion strains, grown from fresh colonies and measured in the same
run, listed on adjacent rows of the released table (each row already the mean of up to
two traces, so the instrument's repeat precision on one extract is not recoverable).
Their rows spread as widely as the whole table (the $\log_{10}$ ratio's sd over the 17
peaks 0.10 to 0.53 against 0.11 to 0.51, medians 0.29 and 0.28), so the low agreement
is not a restricted range: a typical gene's released value is not reproduced by the
same study's second measurement of it, and 0.05 is the reliability ceiling of a Cooper
strain-level value, against M\"ulleder's 0.20 from the same kind of re-measurement.
A cross-screen
correlation cannot exceed roughly the geometric mean of the two, about 0.11, and the
observed 0.00 to 0.06 is at or below it. Cooper's released table carries almost no
strain-level signal, so it can neither validate the M\"ulleder target strain by strain
nor serve as a second one; and the medium and the auxotrophy differ between the screens,
a biological difference that this replicate agreement cannot separate from noise. What
the read fixes is the amino-acid endpoint's position: its own retest rank is 0.20, the
second screen agrees with it at 0.03, and every model above, transformer or ladder,
sits between 0.00 and 0.06. The retest is the agreement of two single-culture
measurements of one strain, so it bounds a second screen, not a model: a model scored
against one measurement is bounded by the square root of the reliability, the share of
a measurement's variance that is between strains (Table~\ref{tab:mulleder-retest}). On
the Pearson scale over $\log_{10}$ concentrations that bound averages 0.62 over the
nineteen amino acids, from 0.00 for asparagine and 0.13 for methionine, whose retest
spread equals their spread across strains, to 0.85 for glutamine; from the Spearman
retest of 0.20 the same rule gives about 0.45. The 237 injections of the quality-control
sample put the instrument's spread at 0.03 to 0.13 in $\log_{10}$ against a
single-culture spread of 0.16 to 0.27, so the retest noise is culture to culture, not
analytical. The models sit at a tenth of that bound, not at it: the target is hard to
learn from single-culture labels, and no input has found more than a tenth of what the
measurement would let a model read.

%% notes-tex/032-amino-acid-module/tables/amino_acid_concordance.tex
%% GENERATED by experiments/031-betaxanthin-module/scripts/target_notes_tex_tables.py
%% SOURCE: results/amino_acid_dataset_concordance.json
\begin{table}[H]
  \centering
  \footnotesize
  \caption[M\"ulleder 2016 against Cooper 2010]{The two deletion-collection amino-acid screens on their 4,051 shared strains, matched on the collection's ORF name. Cooper's CE-LIF peaks (a linear ratio to the plate mean; composite peaks list their components) against M\"ulleder's LC-MS concentrations (mM): the Spearman over strains both measured, the Pearson of the two log-scaled quantities, the overlap of the top 5 percent of strains against the count expected by chance, and the two within-study replicate references, M\"ulleder's (194 strains measured twice, Spearman on log10) and Cooper's duplicated genes (two deletion strains of one gene in the released table). Written by \texttt{target\_notes\_tex\_tables.py} from \texttt{results/amino\_acid\_dataset\_concordance.json}.}
  \label{tab:amino-acid-concordance}
  \begin{tabular}{@{}llrrrrrr@{}}
    \hline
    Cooper peak & M\"ulleder & strains & Spearman & log Pearson & top 5\% & M\"ulleder replicate & Cooper replicate \\
    \hline
    arginine & arginine & 4044 & $-0.000$ & $+0.032$ & 11 / 10.1 & 0.20 & 0.07 \\
    threonine & threonine & 3752 & $+0.007$ & $+0.007$ & 7 / 9.4 & 0.19 & 0.03 \\
    alanine & alanine & 4048 & $-0.006$ & $+0.015$ & 11 / 10.1 & 0.15 & -0.02 \\
    serine & serine & 3756 & $+0.007$ & $+0.022$ & 15 / 9.4 & 0.29 & 0.00 \\
    glycine & glycine & 4044 & $+0.059$ & $+0.047$ & 10 / 10.1 & 0.05 & 0.13 \\
    glutamate & glutamate & 4021 & $+0.043$ & $+0.050$ & 10 / 10.0 & 0.30 & 0.01 \\
    aspartate & aspartate & 3736 & $+0.031$ & $+0.035$ & 11 / 9.4 & 0.22 & 0.04 \\
    lysine a & lysine & 4045 & $+0.021$ & $+0.043$ & 19 / 10.1 & 0.10 & 0.17 \\
    lysine b & lysine & 4032 & $+0.029$ & $+0.057$ & 22 / 10.1 & 0.10 & 0.02 \\
    leucine+isoleucine+citrulline & leucine+isoleucine & 4043 & $+0.019$ & $+0.025$ & 15 / 10.1 & 0.23 & 0.12 \\
    leucine+isoleucine+citrulline & leucine & 4043 & $+0.007$ & $+0.025$ & 13 / 10.1 & 0.24 & 0.12 \\
    leucine+isoleucine+citrulline & isoleucine & 4043 & $+0.028$ & $+0.021$ & 18 / 10.1 & 0.22 & 0.12 \\
    glutamine+valine & glutamine+valine & 4048 & $+0.009$ & $-0.010$ & 12 / 10.1 & 0.28 & 0.10 \\
    glutamine+valine & glutamine & 4048 & $+0.009$ & $-0.015$ & 12 / 10.1 & 0.34 & 0.10 \\
    glutamine+valine & valine & 4048 & $+0.013$ & $+0.035$ & 11 / 10.1 & 0.21 & 0.10 \\
    methionine+proline & methionine+proline & 4037 & $+0.010$ & $+0.032$ & 17 / 10.1 & 0.24 & 0.26 \\
    methionine+proline & methionine & 4037 & $+0.016$ & $+0.022$ & 6 / 10.1 & 0.21 & 0.26 \\
    methionine+proline & proline & 4037 & $+0.006$ & $+0.029$ & 18 / 10.1 & 0.27 & 0.26 \\
    asparagine+tyrosine & asparagine+tyrosine & 4044 & $+0.031$ & $+0.024$ & 8 / 10.1 & 0.12 & -0.12 \\
    asparagine+tyrosine & asparagine & 4044 & $+0.011$ & $+0.006$ & 14 / 10.1 & 0.18 & -0.12 \\
    asparagine+tyrosine & tyrosine & 4044 & $+0.029$ & $+0.023$ & 8 / 10.1 & 0.07 & -0.12 \\
    \hline
  \end{tabular}
\end{table}


%% notes-tex/032-amino-acid-module/tables/mulleder_retest.tex
%% GENERATED by experiments/031-betaxanthin-module/scripts/target_notes_tex_tables.py
%% SOURCE: experiments/027-betaxanthin-metabolic/results/mulleder_retest_reliability.json
\begin{table}[H]
  \centering
  \footnotesize
  \caption[M\"ulleder 2016 retest and the model ceiling]{M\"ulleder 2016's retest, the 194 strains measured twice, per amino acid. Pearson and Spearman between the two $\log_{10}$ measurements; the quality-control sample's sd over its 237 injections (the instrument); the single-measurement sd from the pairs and the sd across strains (both $\log_{10}$); the reliability, the between-strain share of a measurement's variance; and its square root, the Pearson a perfect model of the true value would reach against one measurement. The retest correlation bounds a second screen; the square root of the reliability bounds a model. Written by \texttt{target\_notes\_tex\_tables.py} from the 027 read \texttt{results/mulleder\_retest\_reliability.json}.}
  \label{tab:mulleder-retest}
  \begin{tabular}{@{}lrrrrrrrr@{}}
    \hline
    amino acid & pairs & Pearson & Spearman & QC sd & single sd & across sd & reliability & ceiling \\
    \hline
    alanine & 193 & 0.28 & 0.15 & 0.03 & 0.16 & 0.21 & 0.40 & 0.64 \\
    aspartate & 193 & 0.46 & 0.22 & 0.03 & 0.16 & 0.22 & 0.43 & 0.66 \\
    glutamate & 194 & 0.67 & 0.30 & 0.05 & 0.21 & 0.36 & 0.66 & 0.82 \\
    phenylalanine & 194 & 0.46 & 0.16 & 0.05 & 0.21 & 0.29 & 0.45 & 0.67 \\
    glycine & 190 & 0.12 & 0.05 & 0.04 & 0.19 & 0.21 & 0.20 & 0.45 \\
    histidine & 194 & 0.50 & 0.26 & 0.13 & 0.18 & 0.26 & 0.49 & 0.70 \\
    isoleucine & 194 & 0.52 & 0.22 & 0.05 & 0.21 & 0.30 & 0.52 & 0.72 \\
    lysine & 194 & 0.48 & 0.10 & 0.07 & 0.18 & 0.23 & 0.43 & 0.65 \\
    leucine & 194 & 0.49 & 0.24 & 0.05 & 0.22 & 0.30 & 0.47 & 0.69 \\
    methionine & 189 & 0.19 & 0.21 & 0.05 & 0.27 & 0.27 & 0.02 & 0.13 \\
    asparagine & 190 & 0.37 & 0.18 & 0.04 & 0.22 & 0.22 & 0.00 & 0.00 \\
    proline & 194 & 0.64 & 0.27 & 0.04 & 0.22 & 0.33 & 0.57 & 0.76 \\
    glutamine & 194 & 0.67 & 0.34 & 0.05 & 0.23 & 0.43 & 0.72 & 0.85 \\
    arginine & 194 & 0.37 & 0.20 & 0.12 & 0.26 & 0.36 & 0.50 & 0.71 \\
    serine & 194 & 0.54 & 0.29 & 0.05 & 0.21 & 0.31 & 0.55 & 0.74 \\
    threonine & 194 & 0.56 & 0.19 & 0.03 & 0.21 & 0.31 & 0.55 & 0.74 \\
    valine & 194 & 0.48 & 0.21 & 0.04 & 0.22 & 0.31 & 0.51 & 0.72 \\
    tryptophan & 190 & 0.24 & 0.14 & 0.06 & 0.19 & 0.22 & 0.24 & 0.49 \\
    tyrosine & 193 & 0.20 & 0.07 & 0.13 & 0.25 & 0.30 & 0.32 & 0.57 \\
    \hline
    mean & & 0.43 & 0.20 & & & & 0.42 & 0.62 \\
    \hline
  \end{tabular}
\end{table}


\subsection{The betaxanthin endpoint against the amino-acid pool\secstatus{todo}}
\label{sec:cross-target}

The three screens profile the same deletion collection, so a gene both screens hold is
one paired observation of the two endpoints, and whether a joint model of betaxanthin
and the amino acids has anything to share is a question the data answer before any model
is trained (Figure~\ref{fig:cross-target}, Table~\ref{tab:betaxanthin-amino-acid-correlation};
\texttt{betaxanthin\_amino\_acid\_correlation.py}). Over the 4,477 strains Cachera and
M\"ulleder both measured, the betaxanthin score correlates with methionine at $+0.130$
(bootstrap interval 0.100 to 0.157), with threonine at $+0.093$ and aspartate at
$+0.091$, and against tyrosine at $-0.083$, arginine at $-0.077$ and glutamate at
$-0.065$; the top 5 percent of strains on the two axes overlap at 3.3 times chance for
aspartate (37 of 224 against 11.2) and 3.0 times for threonine, and at chance for the
rest. Each of these sits well under its ceiling, the correlation two perfectly related
quantities would show through the betaxanthin score's reliability of 0.84 and the amino
acid's retest (0.20 to 0.54 over the nineteen), so the single-amino-acid reads are weak
but not at their floor. The pool read is the one that matters for a joint model: a ridge
from the nineteen standardized log concentrations to the betaxanthin score, fit and
scored out of fold over five strain folds, reads a Spearman of $+0.226$ (folds 0.182 to
0.272) against $-0.002$ (sd 0.016) with the target shuffled. A deletion's measured
amino-acid pool carries about as much rank on its betaxanthin level as the best
genotype-side model in the betaxanthin study reads from the deleted gene (0.211,
\texttt{notes-tex/031-betaxanthin-module}, round 151), on a different set of strains and
from a different kind of input, so the two endpoints share strain-level signal well
above the noise. Whether the pool adds to the deleted gene or restates it is read on the
same strains and folds: the same ridge from the gene's propagated GO multi-hot (the
betaxanthin study's best genotype-side input, 2,362 columns, an all-zero row for the
557 strains outside the store) reads 0.165, and from GO with the pool appended, the GO
block scaled to the pool block's variance so one penalty serves both, 0.237 (folds 0.207
to 0.311), above the pool alone on four folds of five and above GO alone on every fold.
The pool adds 0.072 beyond the gene's identity, so the amino-acid signal a joint model
would learn from is not the one the annotation already carries. The Cooper screen gives
the same picture faintly: no single peak beyond $\pm 0.063$ over the 4,225 shared
strains, a pool ridge at $+0.083$ against a shuffled $-0.003$, with the peaks' own
duplicate agreement bounding the ceilings at 0.00 to 0.49; on these strains GO alone
reads 0.193 and GO with the Cooper pool 0.196, an addition of 0.003. Hypothesis
(untested): part of the shared signal is a growth or respiratory covariate common to
both screens, since the betaxanthin screen's top hits are mitochondrial and the pool
responds to respiratory deficiency.

\tcfigfit[1.0]{figures/betaxanthin_amino_acid_correlation_2026-09-24-21-02-01.pdf}{\textbf{The
betaxanthin score against each amino-acid quantity, strain by strain.}
\textbf{a}, Cachera 2023's betaxanthin score against each of M\"ulleder 2016's nineteen
concentrations over the 4,477 strains both screens hold: the Spearman (bars) with a
strain-bootstrap 95 percent interval (1,000 draws) and, above each bar, the ceiling a
correlation could reach through both endpoints' noise, the geometric mean of the
betaxanthin reliability (0.84) and the amino acid's log10 retest Spearman.
\textbf{b}, The same against Cooper 2010's seventeen peaks over the 4,225 shared strains,
the ceiling from Cooper's duplicate-strain agreement (absent where no duplicate pairs
were measured). Written by \texttt{betaxanthin\_amino\_acid\_correlation.py}.}{fig:cross-target}

%% notes-tex/032-amino-acid-module/tables/betaxanthin_amino_acid_correlation.tex
%% GENERATED by experiments/031-betaxanthin-module/scripts/target_notes_tex_tables.py
%% SOURCE: results/betaxanthin_amino_acid_correlation.json
\begin{table}[H]
  \centering
  \footnotesize
  \caption[Betaxanthin against the amino-acid pool]{Cachera 2023's betaxanthin score (corrected 24 h colony fluorescence, population-centered) against each of M\"ulleder 2016's amino-acid concentrations, over the strains both screens hold, matched on the systematic ORF: the Spearman with a strain-bootstrap 95 percent interval (1,000 draws), the overlap of the top 5 percent of strains on the two axes against the count expected by chance, and the ceiling a correlation could reach through both endpoints' noise, the geometric mean of the betaxanthin reliability (0.84, the 027 variance-components read) and the amino acid's log10 retest Spearman over 194 strains. The pool row is a ridge from the standardized log profile to the betaxanthin score over five strain folds, the alpha chosen on an inner split, scored out of fold, with the same fit against a shuffled target. Written by \texttt{target\_notes\_tex\_tables.py} from \texttt{results/betaxanthin\_amino\_acid\_correlation.json} (\texttt{betaxanthin\_amino\_acid\_correlation.py}).}
  \label{tab:betaxanthin-amino-acid-correlation}
  \begin{tabular}{@{}lrrlrr@{}}
    \hline
    M\"ulleder 2016 (mM) & strains & Spearman & 95\% interval & top 5\% & ceiling \\
    \hline
    alanine & 4,477 & $-0.020$ & $[-0.050, +0.010]$ & 24 / 11.2 & 0.35 \\
    arginine & 4,477 & $-0.077$ & $[-0.108, -0.047]$ & 4 / 11.2 & 0.41 \\
    asparagine & 4,477 & $+0.010$ & $[-0.020, +0.038]$ & 23 / 11.2 & 0.39 \\
    aspartate & 4,477 & $+0.091$ & $[+0.061, +0.122]$ & 37 / 11.2 & 0.43 \\
    glutamate & 4,477 & $-0.065$ & $[-0.095, -0.036]$ & 3 / 11.2 & 0.51 \\
    glutamine & 4,477 & $+0.003$ & $[-0.028, +0.034]$ & 5 / 11.2 & 0.54 \\
    glycine & 4,477 & $+0.046$ & $[+0.018, +0.076]$ & 11 / 11.2 & 0.20 \\
    histidine & 4,477 & $-0.020$ & $[-0.051, +0.010]$ & 7 / 11.2 & 0.46 \\
    isoleucine & 4,477 & $-0.055$ & $[-0.086, -0.023]$ & 17 / 11.2 & 0.43 \\
    leucine & 4,477 & $-0.013$ & $[-0.044, +0.018]$ & 6 / 11.2 & 0.45 \\
    lysine & 4,477 & $-0.050$ & $[-0.078, -0.021]$ & 9 / 11.2 & 0.30 \\
    methionine & 4,477 & $+0.130$ & $[+0.100, +0.157]$ & 14 / 11.2 & 0.42 \\
    phenylalanine & 4,477 & $-0.037$ & $[-0.065, -0.009]$ & 7 / 11.2 & 0.36 \\
    proline & 4,477 & $+0.014$ & $[-0.019, +0.044]$ & 7 / 11.2 & 0.48 \\
    serine & 4,477 & $-0.014$ & $[-0.042, +0.016]$ & 4 / 11.2 & 0.50 \\
    threonine & 4,477 & $+0.093$ & $[+0.060, +0.121]$ & 34 / 11.2 & 0.40 \\
    tryptophan & 4,477 & $+0.001$ & $[-0.027, +0.032]$ & 13 / 11.2 & 0.34 \\
    tyrosine & 4,477 & $-0.083$ & $[-0.112, -0.054]$ & 23 / 11.2 & 0.23 \\
    valine & 4,477 & $-0.024$ & $[-0.054, +0.009]$ & 17 / 11.2 & 0.42 \\
    pool ridge, 19 columns, out of fold & 4,477 & $+0.226$ & folds $+0.182$ to $+0.272$ & permuted $-0.002$ (sd 0.016) & \\
    \hline
  \end{tabular}
\end{table}

\begin{table}[H]
  \centering
  \footnotesize
  \caption[Betaxanthin against the Cooper peaks]{Table~\ref{tab:betaxanthin-amino-acid-correlation} for Cooper 2010's peaks (a linear ratio to the plate mean), the ceiling from the peak's duplicate-strain Spearman, blank where no duplicate pairs were measured.}
  \label{tab:betaxanthin-amino-acid-correlation-cooper}
  \begin{tabular}{@{}lrrlrr@{}}
    \hline
    Cooper 2010 (ratio) & strains & Spearman & 95\% interval & top 5\% & ceiling \\
    \hline
    lysine related peak1 & 1,962 & $-0.037$ & $[-0.077, +0.007]$ & 1 / 4.9 &  \\
    arginine & 4,217 & $+0.008$ & $[-0.020, +0.037]$ & 13 / 10.6 & 0.24 \\
    gshbiotin & 2,772 & $+0.049$ & $[+0.014, +0.086]$ & 5 / 7.0 & 0.40 \\
    n acetylornithine & 3,229 & $-0.047$ & $[-0.083, -0.011]$ & 9 / 8.0 & 0.00 \\
    leucine+isoleucine+citrulline & 4,217 & $-0.062$ & $[-0.093, -0.031]$ & 6 / 10.6 & 0.31 \\
    glutamine+valine & 4,222 & $+0.022$ & $[-0.009, +0.051]$ & 13 / 10.5 & 0.29 \\
    methionine+proline & 4,211 & $+0.005$ & $[-0.025, +0.036]$ & 6 / 10.6 & 0.47 \\
    threonine & 3,899 & $-0.007$ & $[-0.038, +0.025]$ & 13 / 9.8 & 0.15 \\
    alanine & 4,221 & $+0.018$ & $[-0.012, +0.047]$ & 8 / 10.5 & 0.00 \\
    serine & 3,910 & $+0.011$ & $[-0.021, +0.042]$ & 11 / 9.8 & 0.04 \\
    asparagine+tyrosine & 4,216 & $-0.003$ & $[-0.033, +0.027]$ & 11 / 10.6 & 0.00 \\
    glycine & 4,217 & $+0.012$ & $[-0.022, +0.043]$ & 13 / 10.6 & 0.34 \\
    lysine a & 4,219 & $-0.057$ & $[-0.090, -0.025]$ & 12 / 10.6 & 0.38 \\
    ornithine & 3,948 & $-0.056$ & $[-0.090, -0.024]$ & 7 / 9.8 & 0.49 \\
    lysine b & 4,207 & $-0.063$ & $[-0.096, -0.031]$ & 14 / 10.5 & 0.13 \\
    glutamate & 4,189 & $-0.043$ & $[-0.074, -0.012]$ & 14 / 10.4 & 0.10 \\
    aspartate & 3,880 & $-0.042$ & $[-0.075, -0.009]$ & 9 / 9.7 & 0.18 \\
    pool ridge, 17 columns, out of fold & 4,225 & $+0.092$ & folds $+0.066$ to $+0.121$ & permuted $-0.006$ (sd 0.014) & \\
    \hline
  \end{tabular}
\end{table}


